\PassOptionsToPackage{table}{xcolor}
\documentclass[10pt,aspectratio=169]{beamer}
\newcommand{\LectureNumber}{10}
\input{course-theme.tex}
\title{While loops}
\subtitle{CSC103 Programming Fundamentals / Lecture 10}
\author{Dr. Muhammad Farhan}
\institute{Associate Professor of Computer Science\\Department of Computer Science\\COMSATS University Islamabad\\Sahiwal Campus}
\date{Fall 2026}
\begin{document}
\begin{frame}\titlepage\end{frame}

\begin{frame}[fragile,t]{The problem for this lecture}
\begin{columns}[T,onlytextwidth]\begin{column}{0.60\textwidth}
Design a sentinel loop for marks ending at -1. Print the count and average, or No data when the sentinel arrives first.
\end{column}\begin{column}{0.35\textwidth}\small
Initialise sum and count to zero, then read the first mark.
\end{column}\end{columns}
\note{Introduce the target problem before explaining the tools. Revisit it in the guided solution later.\par Syllabus reading: Deitel Ch 8 (as listed). CLO-2.}
\end{frame}

\begin{frame}[fragile,t]{Learning objectives}
\begin{columns}[T,onlytextwidth]\begin{column}{0.60\textwidth}
\begin{itemize}
\item Design terminating while loops
\item Use counter, sentinel and flag control
\item Trace state across iterations
\end{itemize}
\end{column}\begin{column}{0.35\textwidth}\small
CLO-2

The while structure\par Counter-controlled repetition\par Sentinel-controlled repetition\par Flag-controlled repetition
\end{column}\end{columns}
\note{Explain how each objective supports the opening problem. The final questions check these objectives.\par Syllabus reading: Deitel Ch 8 (as listed). CLO-2.}
\end{frame}

\begin{frame}[fragile,t]{The while structure}
\begin{itemize}
\item A while loop checks its condition before each iteration.
\item Its body may run zero times.
\item Initialisation, condition and progress must agree.
\end{itemize}
\vspace{1mm}\begin{block}{Example}
\begin{lstlisting}
int i = 1;
while (i <= 3) {
  System.out.println(i);
  i++;
}
\end{lstlisting}
\end{block}
\note{Trace i=1,2,3,4. The last condition check is false and does not execute the body.\par Syllabus reading: Deitel Ch 8 (as listed). CLO-2.}
\end{frame}

\begin{frame}[fragile,t]{Counter-controlled repetition}
\begin{itemize}
\item Use a counter when the number of repetitions is known.
\item An accumulator starts at an appropriate identity, such as zero for a sum.
\item State the range clearly.
\end{itemize}
\vspace{1mm}\begin{block}{Example}
\begin{lstlisting}
int total = 0;
int i = 1;
while (i <= 4) {
  total += i;
  i++;
}
\end{lstlisting}
\end{block}
\note{The sum after each iteration is 1,3,6,10. An invariant is that total contains the sum of all already processed values.\par Syllabus reading: Deitel Ch 8 (as listed). CLO-2.}
\end{frame}

\begin{frame}[fragile,t]{Sentinel-controlled repetition}
\begin{itemize}
\item A sentinel marks the end of input and is not ordinary data.
\item Read before the loop and again after processing each accepted value.
\item Do not include the sentinel in a sum or count.
\end{itemize}
\vspace{1mm}\begin{block}{Example}
\begin{lstlisting}
int x = in.nextInt();
while (x != -1) {
  total += x;
  count++;
  x = in.nextInt();
}
\end{lstlisting}
\end{block}
\note{Assume valid integer tokens and nonnegative data. Test the sentinel as the first input. Check count before calculating an average.\par Syllabus reading: Deitel Ch 8 (as listed). CLO-2.}
\end{frame}

\begin{frame}[fragile,t]{Flag-controlled repetition}
\begin{itemize}
\item A Boolean flag can record whether a condition has been achieved.
\item Every path must permit progress toward stopping.
\item A loop that never changes its condition may not terminate.
\end{itemize}
\vspace{1mm}\begin{block}{Example}
\begin{lstlisting}
boolean found = false;
int candidate = 1;
while (!found && candidate <= 10) {
  found = candidate % 7 == 0;
  candidate++;
}
\end{lstlisting}
\end{block}
\note{This loop stops after finding 7. The increment leaves candidate at 8, illustrating why the final counter and matched value can differ.\par Syllabus reading: Deitel Ch 8 (as listed). CLO-2.}
\end{frame}

\begin{frame}[fragile,t]{A loop maintains a useful relationship}
\begin{itemize}
\item After processing k marks, count should equal k and sum should contain their total.
\item The next iteration extends this relationship with one accepted mark.
\item This makes the final calculation easier to justify.
\end{itemize}
\vspace{1mm}\begin{block}{Example}
\begin{lstlisting}
Before input: count=0, sum=0
After 60: count=1, sum=60
After 80: count=2, sum=140
\end{lstlisting}
\end{block}
\note{Explain the mechanism using the displayed example. After processing k marks, count should equal k and sum should contain their total. The next iteration extends this relationship with one accepted mark. This makes the final calculation easier to justify.\par Syllabus reading: Deitel Ch 8 (as listed). CLO-2.}
\end{frame}

\begin{frame}[fragile,t]{The sentinel is control data}
\begin{itemize}
\item The stop value ends repetition and contributes neither a mark nor a count.
\item Checking before accumulation prevents a biased sum.
\item An immediate sentinel leaves count at zero, so there is no average to divide for.
\end{itemize}
\vspace{1mm}\begin{block}{Example}
\begin{lstlisting}
Input: 60 80 -1
Data values: 60 and 80
Sentinel: -1
\end{lstlisting}
\end{block}
\note{Explain the mechanism using the displayed example. The stop value ends repetition and contributes neither a mark nor a count. Checking before accumulation prevents a biased sum. An immediate sentinel leaves count at zero, so there is no average to divide for.\par Syllabus reading: Deitel Ch 8 (as listed). CLO-2.}
\end{frame}


\begin{frame}[fragile,t]{A while loop checks before every iteration}
\begin{center}\begin{tikzpicture}
\node[box] (a) at (-4,0) {Initialise i = 1};
\node[decision] (b) at (0,0) {i \textless{}= 3?};
\node[alt] (c) at (4,0) {Use i; increment};
\node[hot] (d) at (0,-2.0) {Continue after loop};
\draw[arr] (a)--(b);\draw[arr] (b)--node[lab]{true}(c);\draw[arr] (b)--node[lab]{false}(d);\draw[arr] (c.north)--++(0,.8)-| (b.north);
\end{tikzpicture}\end{center}
\begin{block}{What to notice}The back edge returns to the condition. A false first check gives zero iterations.\end{block}
\note{Trace each labelled connection. The back edge returns to the condition. A false first check gives zero iterations.}
\end{frame}
\begin{frame}[fragile,t]{Key distinctions}
\small\renewcommand{\arraystretch}{1.5}
\rowcolors{2}{pale}{white}
\begin{tabularx}{\textwidth}{>{\raggedright\arraybackslash}X>{\raggedright\arraybackslash}X>{\raggedright\arraybackslash}X}
\rowcolor{navy}
\color{white}\bfseries Control style & \color{white}\bfseries Stop rule & \color{white}\bfseries Typical use\\
Counter & Required number processed & Read exactly 5 marks\\
Sentinel & Special input encountered & Read until -1\\
Flag & State becomes true or false & Stop after a match\\
\end{tabularx}
\vspace{3mm}\footnotesize These distinctions determine which operation or control structure fits the problem.
\note{\par Syllabus reading: Deitel Ch 8 (as listed). CLO-2.}
\end{frame}

\begin{frame}[fragile,t]{First example: execution}
\begin{lstlisting}
int i = 1;
int total = 0;
while (i <= 4) {
  total += i;
  i++;
}
System.out.println(total);
\end{lstlisting}
\note{This is the main body from Java\_Examples/Lecture10.java. The source file includes the complete class. Predict the result before the next slide.\par Syllabus reading: Deitel Ch 8 (as listed). CLO-2.}
\end{frame}

\begin{frame}[fragile,t]{First example: result}
\begin{columns}[T,onlytextwidth]\begin{column}{0.60\textwidth}
10\par After bodies: (i,total) = (2,1), (3,3), (4,6), (5,10).\par The condition fails when i is 5.
\end{column}\begin{column}{0.35\textwidth}\small
Sentinel-controlled repetition

Flag-controlled repetition
\end{column}\end{columns}
\note{Expected output and interpretation: 10
After bodies: (i,total) = (2,1), (3,3), (4,6), (5,10).
The condition fails when i is 5.\par Syllabus reading: Deitel Ch 8 (as listed). CLO-2.}
\end{frame}

\begin{frame}[fragile,t]{Guided practice}
\begin{columns}[T,onlytextwidth]\begin{column}{0.60\textwidth}
Design a sentinel loop for marks ending at -1. Print the count and average, or No data when the sentinel arrives first.
\end{column}\begin{column}{0.35\textwidth}\small
Attempt the solution before continuing.

Use the definitions and examples from this lecture.
\end{column}\end{columns}
\note{Give students 8 minutes to attempt the problem. The following slides provide a complete solution. Original solution guidance: Initialise sum=0 and count=0. Add only non-sentinel marks. Divide (double) sum by count only when count \textgreater{} 0. Inputs 60,80,-1 give count 2 and average 70.\par Syllabus reading: Deitel Ch 8 (as listed). CLO-2.}
\end{frame}

\begin{frame}[fragile,t]{Solution strategy}
\begin{columns}[T,onlytextwidth]\begin{column}{0.60\textwidth}
\begin{itemize}
\item Initialise sum and count to zero, then read the first mark.
\item While the mark is not -1, accumulate it and read the next mark.
\item If count is positive, divide using floating-point arithmetic. Otherwise print No data.
\end{itemize}
\end{column}\begin{column}{0.35\textwidth}\small
Sample console input\par 60 80 -1\par 
\end{column}\end{columns}
\note{Discuss why each step is needed. State the input assumptions before executing the program.\par Syllabus reading: Deitel Ch 8 (as listed). CLO-2.}
\end{frame}

\begin{frame}[fragile,t]{Complete solution (1/2)}
\begin{lstlisting}
public class Solution10 {
  public static void main(String[] args) throws Exception {
    java.util.Scanner in = new java.util.Scanner(System.in);
    int sum = 0;
    int count = 0;
    int mark = in.nextInt();
    while (mark != -1) {
      sum += mark;
\end{lstlisting}
\note{Complete runnable file: Java\_Examples/Solution10.java. Standard input: 60 80 -1
  \par Syllabus reading: Deitel Ch 8 (as listed). CLO-2.}
\end{frame}

\begin{frame}[fragile,t]{Complete solution (2/2)}
\begin{lstlisting}
      count++;
      mark = in.nextInt();
    }
    if (count == 0) System.out.println("No data");
    else System.out.println((double) sum / count);
  }
}
\end{lstlisting}
\note{Complete runnable file: Java\_Examples/Solution10.java. Standard input: 60 80 -1
  \par Syllabus reading: Deitel Ch 8 (as listed). CLO-2.}
\end{frame}

\begin{frame}[fragile,t]{Execution trace}
\small\renewcommand{\arraystretch}{1.5}
\rowcolors{2}{pale}{white}
\begin{tabularx}{\textwidth}{>{\raggedright\arraybackslash}X>{\raggedright\arraybackslash}X>{\raggedright\arraybackslash}X>{\raggedright\arraybackslash}X}
\rowcolor{navy}
\color{white}\bfseries Input & \color{white}\bfseries mark != -1 & \color{white}\bfseries sum after body & \color{white}\bfseries count after body\\
60 & true & 60 & 1\\
80 & true & 140 & 2\\
-1 & false & 140 unchanged & 2 unchanged\\
\end{tabularx}
\vspace{3mm}\footnotesize Follow the changing state in execution order.
\note{Manually derived trace for the complete solution. Initialise sum and count to zero, then read the first mark. While the mark is not -1, accumulate it and read the next mark. If count is positive, divide using floating-point arithmetic. Otherwise print No data.\par Syllabus reading: Deitel Ch 8 (as listed). CLO-2.}
\end{frame}

\begin{frame}[fragile,t]{Solution result}
\begin{columns}[T,onlytextwidth]\begin{column}{0.60\textwidth}
70.0
\end{column}\begin{column}{0.35\textwidth}\small
Why this result follows

If count is positive, divide using floating-point arithmetic. Otherwise print No data.
\end{column}\end{columns}
\note{Expected result: 70.0\par Syllabus reading: Deitel Ch 8 (as listed). CLO-2.}
\end{frame}

\begin{frame}[fragile,t]{A common incorrect approach}
int i = 1;\par while (i \textless{}= 5) \{ System.out.println(i); \}
\vspace{1mm}\begin{block}{Example}
\begin{lstlisting}
i never changes. Add i++ inside the loop so the condition eventually becomes false.
\end{lstlisting}
\end{block}
\note{Ask students to predict the failure before discussing the explanation. The right side gives the correction.\par Syllabus reading: Deitel Ch 8 (as listed). CLO-2.}
\end{frame}

\begin{frame}[fragile,t]{Boundary and failure checks}
\begin{columns}[T,onlytextwidth]\begin{column}{0.60\textwidth}
Check the stated input assumptions.

Use while to compute the digit sum of a nonnegative integer. Test 0, 7 and 204.
\end{column}\begin{column}{0.35\textwidth}\small
Initialise sum=0 and count=0. Add only non-sentinel marks. Divide (double) sum by count only when count \textgreater{} 0. Inputs 60,80,-1 give count 2 and average 70.
\end{column}\end{columns}
\note{Use the specification to derive each expected result. Exercise solution: Initialise sum=0 and count=0. Add only non-sentinel marks. Divide (double) sum by count only when count \textgreater{} 0. Inputs 60,80,-1 give count 2 and average 70.\par Syllabus reading: Deitel Ch 8 (as listed). CLO-2.}
\end{frame}

\begin{frame}[fragile,t]{Check your understanding}
\begin{columns}[T,onlytextwidth]\begin{column}{0.60\textwidth}
Can while run zero times? Why must a sentinel be outside the normal input domain?
\end{column}\begin{column}{0.35\textwidth}\small
Write a prediction and explain the rule that supports it.
\end{column}\end{columns}
\note{Answer: Yes, if the initial condition is false. Otherwise the program cannot distinguish a data value from the stop signal.\par Syllabus reading: Deitel Ch 8 (as listed). CLO-2.}
\end{frame}

\begin{frame}[fragile,t]{Answer and explanation}
\begin{columns}[T,onlytextwidth]\begin{column}{0.60\textwidth}
Yes, if the initial condition is false. Otherwise the program cannot distinguish a data value from the stop signal.
\end{column}\begin{column}{0.35\textwidth}\small
i never changes. Add i++ inside the loop so the condition eventually becomes false.
\end{column}\end{columns}
\note{Reconnect the answer to the relevant definition rather than treating it as a fact to memorise.\par Syllabus reading: Deitel Ch 8 (as listed). CLO-2.}
\end{frame}


\begin{frame}[fragile,t]{Compare cases and predict outcomes}
\small\renewcommand{\arraystretch}{1.5}\rowcolors{2}{pale}{white}
\begin{tabularx}{\textwidth}{>{\raggedright\arraybackslash}X>{\raggedright\arraybackslash}X>{\raggedright\arraybackslash}X}\rowcolor{navy}
\color{white}\bfseries Check & \color{white}\bfseries i & \color{white}\bfseries Action\\
First & 1 & Use 1; set i to 2\\
Third & 3 & Use 3; set i to 4\\
Final & 4 & Exit without using 4\\
\end{tabularx}\par\vspace{3mm}\begin{exampleblock}{Discuss}Explain each expected result before running the example. Which case would reveal a wrong assumption?\end{exampleblock}
\note{Read the first two columns and invite predictions before discussing the outcome.}
\end{frame}

\begin{frame}[fragile,t]{Apply the idea}
\begin{alertblock}{Independent practice}Write a while loop that adds the odd integers 1, 3, 5 and 7. State the final counter value and the number of condition checks.\end{alertblock}\vspace{3mm}\begin{enumerate}\item Work through the input by hand.\item Write or correct the program.\item Justify the expected result.\end{enumerate}\note{Allow 3–5 minutes, then compare answers in pairs. Write a while loop that adds the odd integers 1, 3, 5 and 7. State the final counter value and the number of condition checks.}
\end{frame}

\begin{frame}[fragile,t]{Solution and reasoning}
\begin{lstlisting}
int i = 1;
int sum = 0;
while (i <= 7) {
  sum += i;
  i += 2;
}
System.out.println(sum);
\end{lstlisting}\vspace{1mm}\small\begin{itemize}
\item The four additions give 16.
\item The counter is 9 after the final update.
\item There are five condition checks, including the last false check.
\end{itemize}\note{Answer: The four additions give 16. The counter is 9 after the final update. There are five condition checks, including the last false check.}
\end{frame}

\begin{frame}[fragile,t]{A bounded number{-}guessing loop}
\begin{itemize}
\item The flag records success; the counter limits how long the loop may continue.
\item Increment the attempt count once for each accepted integer guess.
\item Use a fixed secret while tracing so the expected outcome is reproducible.
\end{itemize}
\vspace{3mm}\footnotesize\textcolor{accent}{Source: supplied ISB campus slides. See speaker notes and ISB\_Source\_Map.md.}
\note{Adapted from the supplied ISB campus folder: Lecture{-}10.pptx, slides 12, 13, 14, 22, 23, 24. Replaced randomness with secret 42 for reproducible tracing; numeric input is assumed. Kept the source five{-}attempt limit.}
\end{frame}

\begin{frame}[fragile,t]{Worked problem: A bounded number{-}guessing loop}
\begin{block}{Task}Use secret 42 and at most five guesses. Test input 10, 80, 42 and report Low, High, then Correct.\end{block}
\small Input: 10 80 42\par\vspace{3mm}\begin{exampleblock}{Before running}Predict the output and identify the variables that change.\end{exampleblock}
\note{Adapted from the supplied ISB campus folder: Lecture{-}10.pptx, slides 12, 13, 14, 22, 23, 24. Replaced randomness with secret 42 for reproducible tracing; numeric input is assumed. Kept the source five{-}attempt limit.}
\end{frame}

\begin{frame}[fragile,t]{Complete ISB example (1/2)}
\begin{lstlisting}
public class IsbLecture10 {
  public static void main(String[] args) throws Exception {
    java.util.Scanner in = new java.util.Scanner(System.in);
    int secret = 42, attempts = 0;
    boolean done = false;
    while (attempts < 5 && !done) {
      int guess = in.nextInt();
      attempts++;
      if (guess == secret) {
        done = true;
\end{lstlisting}
\note{Adapted from the supplied ISB campus folder: Lecture{-}10.pptx, slides 12, 13, 14, 22, 23, 24. Replaced randomness with secret 42 for reproducible tracing; numeric input is assumed. Kept the source five{-}attempt limit. Complete file: ISB\_Java\_Examples/IsbLecture10.java.}
\end{frame}

\begin{frame}[fragile,t]{Complete ISB example (2/2)}
\begin{lstlisting}
        System.out.println("Correct");
      } else {
        System.out.println(guess < secret ? "Low" : "High");
      }
    }
    System.out.println("Attempts: " + attempts);
  }

}
\end{lstlisting}
\note{Adapted from the supplied ISB campus folder: Lecture{-}10.pptx, slides 12, 13, 14, 22, 23, 24. Replaced randomness with secret 42 for reproducible tracing; numeric input is assumed. Kept the source five{-}attempt limit. Complete file: ISB\_Java\_Examples/IsbLecture10.java.}
\end{frame}

\begin{frame}[fragile,t]{Worked trace and expected result}
\small\renewcommand{\arraystretch}{1.4}\rowcolors{2}{pale}{white}
\begin{tabularx}{\textwidth}{>{\raggedright\arraybackslash}X>{\raggedright\arraybackslash}X>{\raggedright\arraybackslash}X}\rowcolor{navy}
\color{white}\bfseries Guess & \color{white}\bfseries attempts & \color{white}\bfseries done / message\\
10 & 1 & false / Low\\
80 & 2 & false / High\\
42 & 3 & true / Correct\\
\end{tabularx}\par\vspace{2mm}
\begin{block}{Expected console output}\ttfamily\footnotesize Low\par High\par Correct\par Attempts: 3\end{block}
\note{Adapted from the supplied ISB campus folder: Lecture{-}10.pptx, slides 12, 13, 14, 22, 23, 24. Replaced randomness with secret 42 for reproducible tracing; numeric input is assumed. Kept the source five{-}attempt limit. Trace and expected values independently worked for this edition.}
\end{frame}

\begin{frame}[fragile,t]{Extend the ISB example}
\begin{alertblock}{Practice}Give an input sequence that exercises termination by the attempt limit rather than success.\end{alertblock}\vspace{3mm}\begin{exampleblock}{Explain your reasoning}Five wrong integers, such as 1 2 3 4 5, produce five attempts with done still false.\end{exampleblock}
\note{Adapted from the supplied ISB campus folder: Lecture{-}10.pptx, slides 12, 13, 14, 22, 23, 24. Replaced randomness with secret 42 for reproducible tracing; numeric input is assumed. Kept the source five{-}attempt limit. Answer: Five wrong integers, such as 1 2 3 4 5, produce five attempts with done still false.}
\end{frame}
\begin{frame}[fragile,t]{Lecture summary}
\begin{columns}[T,onlytextwidth]\begin{column}{0.60\textwidth}
\begin{itemize}
\item Design terminating while loops
\item counter, sentinel and flag control
\item state across iterations
\end{itemize}
\end{column}\begin{column}{0.35\textwidth}\small
\begin{itemize}
\item Initialise sum and count to zero, then read the first mark.
\item While the mark is not -1, accumulate it and read the next mark.
\item If count is positive, divide using floating-point arithmetic. Otherwise print No data.
\end{itemize}
\end{column}\end{columns}
\note{Ask students to give one concrete example for the concepts listed. Use the worked solution to connect the topics.\par Syllabus reading: Deitel Ch 8 (as listed). CLO-2.}
\end{frame}

\begin{frame}[fragile,t]{After-class practice}
\begin{columns}[T,onlytextwidth]\begin{column}{0.60\textwidth}
Use while to compute the digit sum of a nonnegative integer. Test 0, 7 and 204.
\end{column}\begin{column}{0.35\textwidth}\small
Syllabus reading\par Deitel Ch 8 (as listed)

Next lecture: For and do-while loops
\end{column}\end{columns}
\note{Reading labels follow the supplied outline. Check topic headings against the available textbook edition. Require a program and expected results for the practice task.\par Syllabus reading: Deitel Ch 8 (as listed). CLO-2.}
\end{frame}
\end{document}
